\chapter{测地线、完备性与变分}
\label{ch:12-geodesics}
\ChapterPlan{周4.2及4.4（本地页码待核对）；梁上册第二版3.3、习题9和10；Meinrenken 16--17及L\"oh 3.5补充}{正文及证明初稿；附8题详解与3幅原创图，尚未独立数学审阅}

第10章给出测地线初值问题，第11章证明Levi--Civita联络使其速度范数恒定。本章研究两类新的问题：给定两端点，测地线是否存在、是否最短；沿一条测地线改变初始数据，邻近测地线如何分离或重新汇合。前者连接长度距离与完备性，后者由Jacobi方程和第二变分刻画。

除非另说，$(M,g)$为有限维、光滑、无边界黎曼流形；讨论全局距离与Hopf--Rinow时另假设连通。联络均为Levi--Civita联络，曲率仍取第10章约定。因此Jacobi方程中为$D_t^2J+R(J,\dot\gamma)\dot\gamma=0$，正曲率球面的法向系数满足正弦型方程。全章不将Lorentz情形的极长性或因果结论混入正定度量论证。

\Needspace{17\baselineskip}
\section{指数映射、法坐标与Gauss引理}
先约定分段光滑曲线$c:[a,b]\to M$的长度为$L(c)=\int_a^b|\dot c|_g\,\dd t$。积分按有限分段相加，允许驻点和常值段；速度不为零时可作弧长重参数化。关于长度距离与能量的详细讨论在下一节，以下局部结论只用此积分。

\begin{definition}[指数映射的自然定义域]\label{def:exponential-domain}
令$\gamma_v$为初值$(p,v)$的极大测地线，$I_v$为其含零的极大时间区间。定义
\[
 \mathcal D=\{(p,v)\in TM:1\in I_v\},\qquad
 \exp_p(v)=\gamma_v(1),\quad v\in\mathcal D_p=\mathcal D\cap T_pM.
\]
不预先假设测地线完备；$\mathcal D_p$一般只是切空间中包含零的开子集。
\end{definition}

\begin{theorem}[指数映射与法坐标]\label{thm:normal-coordinates}
$\mathcal D\subset TM$开且包含零截面，$(p,v)\mapsto\exp_pv$光滑。每个$\mathcal D_p$关于零星形，且在有定义的时间上$\exp_p(tv)=\gamma_v(t)$。微分$(\dd\exp_p)_0$为恒等映射，故$\exp_p$把零附近某开球微分同胚到$p$的邻域。

以$T_pM$的正交归一基识别其为$\R^n$，用$\exp_p^{-1}$得到的坐标称以$p$为中心的黎曼法坐标。在该坐标中
\[
 x^i(\gamma_v(t))=tv^i,\qquad
 g_{ij}(p)=\delta_{ij},\quad \Gamma^k_{ij}(p)=0,\quad \partial_k g_{ij}(p)=0.
\]
第一式只针对从中心出发、在法坐标域内的径向测地线，不说所有测地线都成为直线。
\end{theorem}
\begin{proof}
第10章的测地喷流是$TM$上的光滑向量场。第5章极大局部流的定义域开且流映射光滑，因此取时间1后投影到底空间便得前两项。零初速度的解为常值，故零截面在定义域内。唯一性及仿射重参数化给$\gamma_{av}(t)=\gamma_v(at)$，在两边共同定义的含零区间成立；$0\leq a\leq1$时，$[0,a]\subset I_v$，所以$av\in\mathcal D_p$。该缩放式亦给$\exp_p(tv)=\gamma_v(t)$。

沿切空间直线$sw$在零处求导，
\[
 (\dd\exp_p)_0w=\left.\frac{\dd}{\dd s}\right|_0\exp_p(sw)
 =\dot\gamma_w(0)=w.
\]
逆映射定理给局部微分同胚，缩小为星形球。径向直线的坐标表达随定义而来。原点的度量矩阵等于所选基的Gram矩阵，故为$I$。将$x(t)=tv$代入测地线方程，在$t=0$得$\Gamma^k_{ij}(p)v^iv^j=0$对所有$v$成立。无挠给下槽对称，极化推出每个$\Gamma^k_{ij}(p)=0$；再由$\nabla g=0$得$\partial_k g_{ij}(p)=0$。
\end{proof}

稍后需要基点也变化。映射$\mathcal F(p,v)=(p,\exp_pv)$在每个$(p,0)$的微分可逆：局部丛坐标中为$(a,b)\mapsto(a,a+b)$。所以在对角线$(p,p)$附近，两相近端点有光滑依赖端点的短测地线。这里仅用逆映射定理得到局部结论，不假定全局唯一最短线。

\begin{lemma}[二维参数的导数交换]\label{lem:variation-commutation}
对光滑$F(s,t)$，记$S=\partial_sF$、$T=\partial_tF$，把$D_s,D_t$按沿映射的局部系数定义。则
\[
 D_sT=D_tS,\qquad (D_sD_t-D_tD_s)W=R(S,T)W
\]
对任意沿$F$的光滑向量场$W$成立，即使$F$不是浸入也成立。
\end{lemma}
\begin{proof}
在坐标中，$D_sW^k=\partial_sW^k+\Gamma^k_{ij}(F)(\partial_sF^i)W^j$。对$W=T$与$W=S$相减，混合偏导抵消，剩余$\Gamma^k_{ij}(S^iT^j-T^iS^j)$，由无挠为零。再展开两次导数：含$\partial_sW,\partial_tW$的一阶项相消，$\partial_s\partial_tF$项相消，余项为
\[
 S^iT^jW^m\bigl(\partial_i\Gamma^k_{jm}-\partial_j\Gamma^k_{im}
 +\Gamma^k_{i\ell}\Gamma^\ell_{jm}-\Gamma^k_{j\ell}\Gamma^\ell_{im}\bigr),
\]
正是第10章的$R(S,T)W$。计算按局部截面系数进行，无须将$S,T,W$延成底空间向量场。
\end{proof}

\begin{theorem}[Gauss引理]\label{thm:gauss-lemma}
对$v\in\mathcal D_p$及$w\in T_pM$，用切空间的自然线性识别理解微分的输入，则
\[
 g_{\exp_pv}\bigl((\dd\exp_p)_v v,(\dd\exp_p)_v w\bigr)=g_p(v,w).
\]
在指数映射为微分同胚的星形球内，径向单位向量长为1，并与指数像中的同半径球面正交。
\end{theorem}
\begin{proof}
取$F(s,t)=\exp_p(t(v+sw))$、$0\leq t\leq1$；定义域开及紧时间段的ODE依赖保证小$s$时有定义。各$t$曲线是测地线，速度平方恒为$|v+sw|_p^2$。在$s=0$上，利用交换引理，
\[
 \partial_t g(T,S)=g(T,D_tS)=g(T,D_sT)
 =\tfrac12\partial_s|T|^2=g_p(v,w).
\]
在$t=0$，$S=0$，因此$g(T,S)=t\,g_p(v,w)$。取$t=1$，$T=(\dd\exp_p)_v v$、$S=(\dd\exp_p)_v w$，即得公式。令$w=v$得径向长度；令$w\perp v$得球面正交。$v=0$时也由零点微分恒等直接成立。
\end{proof}

\begin{figure}[htbp]
\centering\DGDiagram{gauss-radial}

\caption{左侧在切空间中以同半径圆标出$v$，输入向量$v,w$视为附着在该点；右侧在$q=\exp_pv$的切空间画其微分像，图中$\dd\exp_p$均在$v$处取值。径向长度保持、径向与角向正交，但角向长度可变，所以整个指数映射不必等距。}
\label{fig:gauss-radial}
\end{figure}

\Needspace{17\baselineskip}
\begin{proposition}[径向测地线的最短性]\label{prop:radial-minimizing}
取$\rho>0$使$\exp_p$在包含$\overline B_{2\rho}(0)$的开集上为微分同胚，记$U_\rho=\exp_p B_\rho(0)$。若$q=\exp_pv$、$|v|<\rho$，则径向段的长度为$|v|$，并在所有连接$p,q$的分段光滑曲线中最短。除去保持方向的重参数化及常值停顿，这条最短曲线唯一。
\end{proposition}
\begin{proof}
在$U_\rho\setminus\{p\}$，把曲线写成$\alpha(t)=\exp_p(r(t)u(t))$，$|u(t)|=1$。Gauss引理给
\[
 |\dot\alpha|^2=(r')^2+
 \bigl|(\dd\exp_p)_{ru}(r\dot u)\bigr|^2\geq(r')^2.
\]
对所有避开$p$的分段积分得$L\geq\int|r'|\geq r(\text{终})-r(\text{始})$。若曲线反复经过$p$，对$r>\varepsilon$的各子段先用此式，再令$\varepsilon\downarrow0$，仍得从$p$到$q$的长度至少$|v|$。等号要求角向导数为零且径向不倒退，因微分可逆，像必沿同一射线单调前进。

还要排除离开法邻域的全局捷径。任意离开$U_\rho$的曲线，第一次到达边界时到达某个$\exp_p(\rho u)$；闭球在更大微分同胚域内，故这个边界描述有效。前段长度至少$\rho>|v|$。因此更短或等长竞争曲线不可能离开$U_\rho$，前述等号分析适用。径向段自身速度常为$|v|$（参数取$[0,1]$），长度正为$|v|$。
\end{proof}

“局部最短”在这里有明确半径与竞争曲线范围。它不推出一条测地线无论延长多久都最短，也不推出任意两个远点之间有最短线。


\Needspace{17\baselineskip}
\section{长度与能量的第一变分}
\begin{definition}[长度距离与能量]\label{def:riemannian-distance}
在连通$M$上，定义$d(p,q)=\inf L(c)$，下确界遍历从$p$到$q$的分段光滑曲线。对固定时间区间$[a,b]$、$a<b$，定义
\[
 E(c)=\frac12\int_a^b|\dot c|^2\,\dd t.
\]
长度与能量是两个泛函。长度在正则重参数化下不变，能量则依赖时间分配。
\end{definition}

\begin{proposition}[距离、拓扑与能量比较]\label{prop:distance-topology}
$d$有限且是距离，其拓扑等于流形原拓扑。对任意固定区间上的分段光滑曲线，
\[
 L(c)^2\leq2(b-a)E(c),
\]
等号当且仅当速度范数几乎处处为常数。因此同区间上的常速长度最短线也使能量最小。
\end{proposition}
\begin{proof}
从$p$可由分段光滑路径到达的点组成开集：到达一点后可在坐标球中接直线。其补集也开，否则一条短坐标线可把不可达点接到可达点。连通性使此集为整个$M$，所以$d<\infty$。反向与拼接给对称、三角不等式，非负显然。

固定$p$的坐标球$\overline B_r(0)$紧含于一张卡，$x(p)=0$。正定性和紧致性给常数$0<\lambda\leq\Lambda$，在闭球内有
\[
 \lambda|\xi|^2\leq g_{ij}(x)\xi^i\xi^j\leq\Lambda|\xi|^2.
\]
任意从$p$离开此球的曲线，在第一次到球面前长度至少$\sqrt\lambda r$；留在球内到$x(q)$的曲线长度至少$\sqrt\lambda|x(q)|$。当$|x(q)|<r$，两类曲线合并给$d(p,q)\geq\sqrt\lambda|x(q)|$，而坐标直线给$d(p,q)\leq\sqrt\Lambda|x(q)|$。球外点的距离至少$\sqrt\lambda r$，所以距离非退化，且两种邻域系统相同。由此每点也有包含在紧坐标闭球内的小闭距离球；此结论还不说大闭球紧。

最后，对函数$|\dot c|$与常数1用Cauchy--Schwarz得到能量不等式，等号为该函数几乎处处恒定。常速最短线$c$满足$E(c)=L(c)^2/(2(b-a))$，而任意竞争曲线$\eta$满足$E(\eta)\geq L(\eta)^2/(2(b-a))\geq E(c)$。
\end{proof}

\begin{definition}[变分与变分场]
光滑变分为$F:(-\varepsilon,\varepsilon)\times[a,b]\to M$，$F(0,t)=\gamma(t)$；记$V(t)=\partial_sF(0,t)$。若$F(s,a)=\gamma(a)$且$F(s,b)=\gamma(b)$，称固定端点变分。此时$V(a)=V(b)=0$。
\end{definition}

每个沿$\gamma$的光滑$V$都能实现为变分场：取$F(s,t)=\exp_{\gamma(t)}(sV(t))$。零截面的指数定义域开、$[a,b]$紧，故可统一选小$\varepsilon$。零点微分为恒等给$\partial_sF(0,t)=V(t)$；若端点$V=0$，该变分确实固定端点。只需有限分段光滑时也可逐段应用。

\Needspace{17\baselineskip}
\begin{theorem}[能量与长度的第一变分]\label{thm:first-variation}
对光滑变分$F$，记$T=\dot\gamma$、$V=\partial_sF|_{s=0}$。则
\[
 \left.\frac{\dd}{\dd s}\right|_0E(F(s,\cdot))
 =[g(V,T)]_a^b-\int_a^b g(V,D_tT)\,\dd t.
\]
若$\gamma$正则、$v=|T|>0$、$U=T/v$，还满足
\[
 \left.\frac{\dd}{\dd s}\right|_0L(F(s,\cdot))
 =[g(V,U)]_a^b-\int_a^b g(V,D_tU)\,\dd t.
\]
因此固定端点下能量驻定当且仅当为仿射测地线；正则曲线的长度驻定当且仅当弧长重参数化后为测地线。
\end{theorem}
\begin{proof}
度量相容和$D_s\partial_tF=D_t\partial_sF$给
\[
 E'(0)=\int_a^b g(D_tV,T)\,\dd t.
\]
对$g(V,T)$积分分部即得第一式。速度在紧区间上有正下界时，足够小的变分仍正则，可对平方根求导，得到$L'(0)=\int g(D_tV,U)$，再积分分部得第二式。

固定端点使边界项为零。任意紧支场均可实现为变分场，取$V=\chi D_tT$，其中$\chi\geq0$为任意内部紧支函数，得到能量驻定必使$D_tT=0$；反向代入即得。长度情形同理用$V=\chi D_tU$。由于$D_tU=vD_\ell U$，$\ell$为弧长，$D_tU=0$等价于单位速度参数下的测地线方程。非仿射重参数化不会改变长度，却可能改变能量驻定性。
\end{proof}

\Needspace{17\baselineskip}
\begin{lemma}[最短折线不能有角]\label{lem:no-geodesic-corner}
两条单位速度光滑曲线在连接点相接，若连接处左右速度$u,v$不同，则对连接点两侧足够短的弧段，可以保持两端点并缩短总长度。特别地，一条分段光滑、分段正则的长度最短线，经弧长参数化后是光滑测地线。
\end{lemma}
\begin{proof}
取连接点$p$的坐标使$g_p$为欧氏内积。两侧距离连接时间$\delta$的坐标为$-\delta u+O(\delta^2)$与$\delta v+O(\delta^2)$。两原弧的总长为$2\delta$，坐标直线连接它们的长度为
\[
 \delta|u+v|_{g_p}+O(\delta^2).
\]
因为$u,v$均单位且不同，$|u+v|<2$，所以小$\delta$时严格缩短；若$u+v=0$，直线长度为$O(\delta^2)$，仍成立。该直线位于小卡中；若需光滑连接，可在两端任意小邻域磨光，长度改变可任意小，严格差额仍保留。

长度最短线在每个光滑正则段内由第一变分得单位速度测地线，连接处不能有角，所以左右速度相同；测地线初值唯一性使相邻段是同一光滑解。常值停顿可先删去。更一般，若一条最短线有驻点，以其中一点为中心取上述法球，足够短子弧留在球内，径向最短性及等号唯一性使其像等于径向线；故在像的意义下仍是同一光滑测地线，不要求原有驻点参数本身正则。
\end{proof}

第一变分判别的是驻定，局部最短性由Gauss引理另证。一般测地线延长后可能失去最短性；长度距离的下确界在不完备流形上甚至可能不被任何曲线达到。


\Needspace{17\baselineskip}
\section{Hopf--Rinow定理}
\begin{definition}[两种完备性]
$(M,d)$为完备距离空间，是指每个$d$-Cauchy序列在$M$中收敛。黎曼流形测地完备，是指每条极大仿射测地线的参数区间都是$\R$。后者等价于所有基点的指数映射都定义在整个切空间上：正向时间可由初速度缩放转到时间1，负向再反转初速度。
\end{definition}

\begin{lemma}[紧小球面上的等号点]\label{lem:hopf-rinow-step}
设$r=d(p,q)>0$。取$0<\varepsilon<r$，足够小使$\exp_p$在更大闭球附近为微分同胚。则存在单位$u\in T_pM$，$x=\exp_p(\varepsilon u)$满足
\[
 d(p,x)=\varepsilon,\qquad d(x,q)=r-\varepsilon.
\]
此引理不要求完备，也未假定$p,q$间已有最短线。
\end{lemma}
\begin{proof}
取$p$到$q$的近似最短曲线$c_j$，$L(c_j)\leq r+1/j$。$q$不在$U_\varepsilon=\exp_pB_\varepsilon$内，否则径向最短性给$d(p,q)<\varepsilon$。每条曲线第一次离开此球时到达$x_j\in\exp_pS_\varepsilon$，前段长度至少$\varepsilon$，所以
\[
 r\leq\varepsilon+d(x_j,q)\leq L(c_j)\leq r+1/j.
\]
球面$\exp_pS_\varepsilon$是紧集，取子列$x_j\to x$。距离函数由三角不等式为1-Lipschitz，所以$d(x,q)=r-\varepsilon$；球面点的径向距离为$\varepsilon$，得到所需结论。
\end{proof}

\Needspace{17\baselineskip}
\begin{theorem}[Hopf--Rinow]\label{thm:hopf-rinow}
连通、有限维、无边界黎曼流形上，下列条件等价：
\begin{enumerate}[label=(\roman*),nosep]
\item 长度距离$d$完备；
\item 每个闭且$d$有界的子集都紧；
\item 测地完备；
\item 存在一个$p$使$\exp_p$定义在整个$T_pM$。
\end{enumerate}
在这些条件下，任意两点之间存在长度为其距离的最短测地线，且
\[
 \overline B_R(p)=\exp_p\{v\in T_pM:|v|\leq R\}.
\]
结论保证存在，不保证最短线唯一。
\end{theorem}
\begin{proof}
\textbf{(i)$\Rightarrow$(iii)。}若某非常值极大测地线有有限右端$b$，其速度范数为常数$c>0$，所以$d(\gamma(t),\gamma(s))\leq c|t-s|$。完备性给$t\uparrow b$时$\gamma(t)\to q\in M$。具体地，先取任意$t_j\uparrow b$得Cauchy极限，再用同一估计得整个曲线趋向$q$。

在$q$附近选紧含于一张卡的闭坐标球$K$，晚期曲线均在其内部。度量下界$g\geq\lambda I$使坐标速度$|\dot x|\leq c/\sqrt\lambda$，因此晚期相空间点$(x,\dot x)$位于$TM$的一个紧集。沿$t_j\uparrow b$取相空间收敛子列；在其极限附近，光滑喷流有统一存在时间$\delta>0$。足够晚的子列点满足$b-t_j<\delta/2$，从该点出发的解便越过$b$，由唯一性与原解相接，矛盾。有限左端同理。常值解本来就完备。

\textbf{(iii)$\Rightarrow$(iv)。}任取$p$即可。接下来只假设(iv)，先证明从这个$p$出发能到达任意$q$的最短线，这是证明的主要步骤。

设$r=d(p,q)>0$，用上一引理取$\varepsilon,u,x$，令$\gamma(t)=\exp_p(tu)$，$0\leq t\leq r$。它因(iv)有定义。起点的等号为$d(\gamma(\varepsilon),q)=r-\varepsilon$。设已有某个$t<r$满足
\[
 d(\gamma(t),q)=r-t.
\]
三角不等式及$\gamma$前段长度$t$立刻给$d(p,\gamma(t))=t$，所以前段已最短。在$\gamma(t)$处再次用小球面引理，取$0<\delta<r-t$及径向单位速度段$\eta$，其终点$y$满足$d(y,q)=r-t-\delta$。于是
\[
 r\leq d(p,y)+d(y,q)\leq t+\delta+r-t-\delta=r.
\]
所有不等式均为等号，拼接$\gamma|_{[0,t]}$与$\eta$得到$p$到$y$的最短线。无角引理使接头处速度相同，初值唯一性给$\eta(s)=\gamma(t+s)$。对$0\leq s\leq\delta$，再用三角不等式及这条拼接线的子段最短性，得到$d(\gamma(t+s),q)=r-t-s$。故等号能从$t$继续到一个右邻域。

令$t_*$为所有满足在$[\varepsilon,t]$上等号恒成立的$t$的上确界。距离连续性使等号在$t_*$仍成立；若$t_*<r$，上一步还能延伸，矛盾。因此$t_*=r$，$d(\gamma(r),q)=0$，从而$\gamma(r)=q$且全段最短。$q=p$取常值段。

\textbf{(iv)$\Rightarrow$(ii)。}刚证明从固定$p$出发的最短线存在，所以每个$d(p,q)\leq R$的点都为某个$|v|\leq R$的指数像；反向由径向段长度得$d(p,\exp_pv)\leq|v|$。因此闭距离球恰为切空间闭球的连续像。有限维闭球紧，所以闭距离球紧；任意闭有界集包含在某个这样的闭球中，因而也紧。

\textbf{(ii)$\Rightarrow$(i)。}Cauchy序列有界，包含在一个紧闭球中，故有收敛子列；Cauchy性使整个序列趋于同一极限。等价性完成后，(i)又给所有基点的测地完备，可对任意新基点重复上面的径向延伸证明，得到任意两点之间的最短线及各基点的闭球公式。
\end{proof}

\begin{figure}[htbp]
\centering\DGDiagram{hopf-rinow-step}

\caption{Hopf--Rinow证明中的距离等号，而非预先画出一条已知最短线：近似最短路径先碰小球面，取紧子列得到$x$；再从已知最短径向段的端点向前取$y$。全程等号迫使接头无角，从而继续同一测地线。虚线只示意距离余量。}
\label{fig:hopf-rinow-step}
\end{figure}

\begin{corollary}[紧流形与有界性]\label{cor:compact-riemannian-complete}
紧的无边界黎曼流形逐连通分支测地完备；连通、完备且有界的黎曼流形紧。
\end{corollary}
\begin{proof}
紧流形的每个连通分支闭且紧；在该分支上$d$诱导原拓扑，所以紧距离空间完备，适用Hopf--Rinow。第二项对整个流形用“闭有界集紧”即可。
\end{proof}

有限维、正定和无边界在证明中均有实际作用。带边区间中测地线可撞到边界，不定度量也没有这里的正长度距离；不能原封不动搬用此定理。完备性还不保证指数映射单射，圆周和球面将给出例子。


\Needspace{17\baselineskip}
\section{Jacobi场、共轭点与第二变分}
\begin{definition}[Jacobi场]\label{def:jacobi-field}
沿仿射测地线$\gamma$，满足
\[
 D_t^2J+R(J,T)T=0,\qquad T=\dot\gamma,
\]
的向量场称Jacobi场。这里$R(J,T)T$的第一个输入是$J$，交换前两个输入会改变方程的号。
\end{definition}

\begin{theorem}[测地变分与Jacobi初值]\label{thm:jacobi-variation}
测地线族$F(s,t)$的变分场是Jacobi场。反之，沿一条紧参数段的任意Jacobi场都可由某个光滑测地线族实现。给定任意$A,B\in T_{\gamma(0)}M$，存在唯一Jacobi场满足$J(0)=A,D_tJ(0)=B$；其解空间维数为$2n$。
\end{theorem}
\begin{proof}
各条$t$曲线为仿射测地线，所以$D_tT=0$。对$s$求导，并用引理\ref{lem:variation-commutation}，
\[
 0=D_sD_tT=D_tD_sT+R(S,T)T=D_t^2S+R(S,T)T.
\]
在$s=0$限制即得Jacobi方程。

沿$\gamma$选正交归一平行标架，把$J$的系数写成列向量$j$，方程为$j''+C(t)j=0$，其中$C$为光滑矩阵。这是线性系统，与第10章平行移动的延拓论证相同，在给定紧时间段上不会爆破：把$(j,j')$写成一阶线性系统，用Gronwall界延拓。因此任意两组初值有唯一解，初值空间为$2n$维。

为实现该解，取$\alpha(s)$使$\alpha(0)=\gamma(0)$、$\alpha'(0)=A$，沿其用平行移动定义
\[
 v(s)=P^\alpha_{s\leftarrow0}\bigl(T(0)+sB\bigr),\qquad
 F(s,t)=\exp_{\alpha(s)}(tv(s)).
\]
有$v(0)=T(0)$及$D_sv(0)=B$。紧时间段上的ODE参数依赖给一个统一小$s$区间；这里原测地线在包含该闭段的开区间上有定义。所得变分场的初值为$A$，由导数交换其初始协变导数为$B$。它已经满足Jacobi方程，唯一性使它等于指定的$J$。
\end{proof}

若$\gamma$为单位速度，则$f(t)=g(J,T)$满足$f''=0$，因为$g(R(J,T)T,T)=0$。所以切向Jacobi场为$(a+bt)T$；若一个Jacobi场在两个不同时刻为零，则其切向分量处处为零。这一点不要求测地线无自交。

\Needspace{17\baselineskip}
\begin{proposition}[指数微分与共轭点]\label{prop:conjugate-exponential}
设$\gamma(t)=\exp_p(tv)$、$|v|=1$、$c>0$。若存在非零Jacobi场满足$J(0)=J(c)=0$，称参数时刻$c$与0沿$\gamma$共轭，重数为这样的Jacobi场空间维数。则它等价于$(\dd\exp_p)_{cv}$奇异，重数等于该微分的核维数，且不超过$n-1$。这是沿指定测地线、指定参数时刻的性质，不只由端点集合决定。
\end{proposition}
\begin{proof}
取初值$J_w(0)=0,D_tJ_w(0)=w$，由变分$\exp_p(t(v+sw))$可得
\[
 J_w(t)=(\dd\exp_p)_{tv}(tw),\qquad
 (\dd\exp_p)_{cv}(w)=\frac1cJ_w(c).
\]
所有在0为零的Jacobi场都由唯一的$w$得到。故核与两端为零的解空间同构；非零初始导数给非零场。微分在有限同维空间中可逆当且仅当核为零，所以等价于局部微分同胚。两端为零的场必法向，其初始导数也垂直于$v$，因此重数至多$n-1$。
\end{proof}

指数映射整体不单射与某点微分奇异是不同现象。圆周上指数映射绕圈不单射，却没有共轭点；球面上对跖点则使角向微分退化。法坐标范围内指数微分可逆，因此没有从中心出发的共轭时刻。

\subsection*{第二变分与指标型}
以下令$\gamma:[0,b]\to M$为单位速度测地线，$b>0$。对沿$\gamma$的连续、分段光滑场$V,W$定义
\[
 I(V,W)=\int_0^b\bigl(g(D_tV,D_tW)-g(R(V,T)T,W)\bigr)\,\dd t.
\]
称$I$为指标型（index form）；它不是整数值的Morse指标。由第11章曲率对称性，$g(R(V,T)T,W)=g(R(W,T)T,V)$，所以$I$为对称双线性型。其在$V(0)=V(b)=0$场上的符号控制固定端点的二阶变化。

\Needspace{17\baselineskip}
\begin{theorem}[固定端点的第二变分]\label{thm:second-variation}
对$\gamma$的光滑固定端点变分，变分场为$V$，有
\[
 E''(0)=I(V,V),\qquad L''(0)=I(V^\perp,V^\perp),
 \quad V^\perp=V-g(V,T)T.
\]
因此最短测地线必对所有光滑端点零场满足$I(V,V)\geq0$；若找到$I(V,V)<0$，则存在相同端点的邻近曲线更短。
\end{theorem}
\begin{proof}
由第一变分，固定端点时$E'(s)=-\int g(S,D_tT)$。在$s=0$原曲线为测地线，故再次求导时含$D_sS$的项为零，而
\[
 D_sD_tT=D_t^2S+R(S,T)T.
\]
所以$E''(0)=-\int g(V,D_t^2V+R(V,T)T)$。积分分部且$V$两端为零，得到$I(V,V)$。

令$h(s,t)=|\partial_tF|^2$，在$s=0$有$h=1$。对$\sqrt h$求两次导数，得到
\[
 L''(0)=E''(0)-\int_0^b g(D_tV,T)^2\,\dd t.
\]
写$V=fT+V^\perp$。因$D_tT=0$，$|D_tV|^2=(f')^2+|D_tV^\perp|^2$，且$g(D_tV,T)=f'$；曲率项只含$V^\perp$。两份$(f')^2$抵消，得长度公式。

若$\gamma$最短，它在固定时间区间上也是能量最小，故$E''(0)\geq0$。反之，任意光滑端点零场都可用沿曲线指数实现为变分场；若$I(V,V)<0$，则$E'(0)=0$、小非零$s$时$E(F_s)<E(\gamma)=b/2$。能量比较给$L(F_s)^2\leq2bE(F_s)<b^2$，所以更短。由此无需再额外假设各条竞争曲线常速。
\end{proof}

对光滑场积分分部还给
\[
 I(V,W)=[g(D_tV,W)]_0^b
 -\int_0^b g(D_t^2V+R(V,T)T,W)\,\dd t.
\]
若$V$仅分段光滑，在分割点$t_j$还须加
$g(D_tV(t_j^-)-D_tV(t_j^+),W(t_j))$。这些接头项不能省略，它们正是内点共轭时刻破坏最短性的关键。

\Needspace{17\baselineskip}
\begin{theorem}[无共轭时的指标引理]\label{thm:index-lemma}
若$(0,b]$内无与0共轭的时刻，则$I$在端点为零的非零连续分段光滑场上正定。更一般，给定$V(0)=0$及终值$V(b)$，存在唯一具有相同端值的Jacobi场$J$，且
\[
 I(V,V)\geq I(J,J),
\]
等号当且仅当$V=J$。此结论是指标型的极小性质，不单独宣称全局最短线唯一。
\end{theorem}
\begin{proof}
沿$\gamma$用正交平行标架，记$C(t)w=R(w,T)T$的矩阵；由对称性$C^{\mathsf T}=C$。取矩阵Jacobi解
\[
 A''+CA=0,\qquad A(0)=0,\quad A'(0)=I.
\]
无共轭条件说明$A(t)$在$0<t\leq b$可逆。Wronskian满足
\[
 (A^{\mathsf T}A'-A'^{\mathsf T}A)'
 =A^{\mathsf T}A''-A''^{\mathsf T}A=0,
\]
初始值为零，故$A^{\mathsf T}A'$对称。

先取端点零的场，系数记$v$。在$t>0$写$v=Af$。还需检查$t=0$不是一个被漏掉的奇点：由微积分基本定理，$A(t)=tB(t)$，$B(t)=\int_0^1A'(st)\,\dd s$光滑且$B(0)=I$；同样$v(t)=tz(t)$，$z(0)=v'(0)$。所以$f=B^{-1}z$在0附近光滑延拓，分段情形逐段成立。

展开$v'=A'f+Af'$，利用$A''=-CA$及Wronskian对称性，逐项整理为
\[
 |v'|^2-\langle Cv,v\rangle
 =|Af'|^2+\frac{\dd}{\dd t}\langle Af,A'f\rangle.
\]
积分后的边界项在0为零（$Af=v=0$且$A'f$有界），在$b$为零（$v(b)=0$）；中间接头因$f$连续而相消。因此$I(V,V)=\int_0^b|Af'|^2\geq0$。等号给$t>0$各段$f'=0$，连续性使$f$恒定，而$A(b)$可逆及$v(b)=0$给$f(b)=0$，所以$V=0$。

对一般终值，$J$的系数为$A(t)A(b)^{-1}v(b)$，存在唯一。令$W=V-J$，其两端为零；Jacobi方程及积分分部给$I(J,W)=0$，所以$I(V,V)=I(J,J)+I(W,W)$，前述正定性给结论及等号条件。
\end{proof}

该证明在固定紧段上使用有限维线性ODE及分段积分，没有预设路径空间的Hilbert流形理论。它也未用“无限维Hessian正定”直接推出函数在某个未定义拓扑中严格极小。

\begin{figure}[htbp]
\centering\DGDiagram{jacobi-models}

\caption{单位速度常截面曲率模型中，法向Jacobi场$J(t)=s_K(t)E(t)$，初值$J(0)=0,D_tJ(0)=E(0)$。两图横轴均为$[0,\pi]$，右图纵轴单独缩放。正曲率系数在$\pi$归零，平坦与负曲率模型不归零。纵轴画系数，不是流形中的实际轨迹。}
\label{fig:jacobi-models}
\end{figure}

\Needspace{17\baselineskip}
\begin{theorem}[内点共轭之后不再最短]\label{thm:conjugate-nonminimizing}
若$0<c<b$且$c$与0沿$\gamma$共轭，则存在光滑端点零场$Z$使$I(Z,Z)<0$，因而$\gamma|_{[0,b]}$不是最短线。若只在终点$b$共轭，以上结论不能直接套用；球面到对跖点仍可最短。
\end{theorem}
\begin{proof}
取非零Jacobi场$J$，$J(0)=J(c)=0$，并将它在$[c,b]$延为零，得到连续分段光滑场$V$。Jacobi积分分部给$I(V,V)=0$。取光滑端点零场$W$，使$W(c)=-D_tJ(c)$；可将$-D_tJ(c)$平行延拓后乘支于$(0,b)$、在$c$为1的光滑截断函数。唯一性使$D_tJ(c)\ne0$，否则在$c$值和导数同为零会使$J$恒零。接头项给
\[
 I(V,W)=g(D_tJ(c),W(c))=-|D_tJ(c)|^2<0.
\]
因此对足够小的$\eta>0$，
\[
 I(V+\eta W,V+\eta W)
 =-2\eta|D_tJ(c)|^2+\eta^2I(W,W)<0.
\]
还须将这个有导数折角的场变成光滑场。在平行标架中，只在$c$附近长$O(\delta)$的小区间内平滑其系数，区间外保持原场及端点不变。具体可先将连续分段光滑系数作卷积，再用等于1的局部截断与原场拼接；系数误差为$O(\delta)$，导数误差在$L^2$中趋零，因为原导数分段有界，拼接产生的截断导数项也只在长$O(\delta)$区间内有界。故所得光滑$Z_\delta$一致趋于原场，且一阶导数$L^2$趋近。

紧时间段上$C(t)$有界，Cauchy--Schwarz说明$I$对这种收敛连续，所以小$\delta$时仍有$I(Z_\delta,Z_\delta)<0$。第二变分定理于是构造真正光滑、固定端点的更短曲线。若$c=b$，所需$W(b)=-D_tJ(b)\ne0$与固定端点条件矛盾，这也解释为什么证明不能覆盖终点共轭情形。
\end{proof}

本章得到的是共轭与局部二阶变化的关系。割点还可能来自两条不同最短线相遇，未必是共轭点。一般割迹结构、Morse指标定理、Rauch比较、Bonnet--Myers及Cartan--Hadamard定理尚未证明；以下只在具体模型中计算相应现象。


\Needspace{15\baselineskip}
\section{测地线与变分习题及详细解答}
习题12.1改述梁书上册第二版第3章原题9，12.2改述原题10，两题均见印刷页85（PDF第101页）。

12.3把正文定理3-3-8（印刷75/PDF91）的径向坐标结论练习化并增加说明，不另造原题号。12.4–12.8为原创补充。全部解答由本工程按既定联络与曲率符号重写。

\Needspace{12\baselineskip}
\subsection*{习题12.1：非仿射参数与常速}
改述原题9。设正则曲线满足$D_t\dot c=\alpha(t)\dot c$，构造使其成为仿射测地线的参数。证明非恒定测地线的两个仿射参数只差$at+b$，并在正定度量下证明弧长参数为仿射参数。
\begin{proof}[解答]
令新参数$s=s(t)$，$s'>0$。链式法则给
\[
 D_s\frac{\dd c}{\dd s}
 =\frac1{(s')^2}\left(D_t\dot c-\frac{s''}{s'}\dot c\right).
\]
取$s''/s'=\alpha$，即在任意给定区间内选
\[
 s'(t)=C\exp\left(\int_{t_0}^t\alpha(u)\,\dd u\right)>0,
 \qquad s(t)=s_0+\int_{t_0}^t s'(u)\,\dd u,
\]
得到到其像区间的光滑微分同胚，方程变为零加速度。这不保证像区间为整个$\R$。若$t$本来仿射，则保持零加速度要求$s''=0$，正则性使$s'=a\ne0$，所以$s=at+b$，允许$a<0$时还可反向。

度量相容给$\partial_t|\dot\gamma|^2=2g(D_t\dot\gamma,\dot\gamma)=0$，非恒定测地线速度为某个$c>0$。弧长$\ell=ct+\ell_0$与$t$为仿射关系，所以也是仿射参数。常值曲线没有可逆弧长参数；Lorentz零测地线的“长度参数”也不能按此构造，本题仅讨论正定度量的正则曲线。
\end{proof}

\Needspace{12\baselineskip}
\subsection*{习题12.2：球面方程与大圆}
改述原题10。在半径$a$球面的坐标度量$a^2(\dd\theta^2+\sin^2\theta\,\dd\phi^2)$下写出测地线方程，并验证配匀速参数的大圆为测地线。说明反向分类为何也成立。
\begin{proof}[解答]
第11章已求系数，代入得
\[
 \ddot\theta-\sin\theta\cos\theta\,\dot\phi^2=0,
 \qquad \ddot\phi+2\cot\theta\,\dot\theta\dot\phi=0.
\]
赤道$\theta=\pi/2$、$\phi=ct+\phi_0$满足两式；绕球心的正交变换保持诱导度量及Levi--Civita联络，所以每条大圆的匀速参数都是测地线。给定$p$及非零$v\in T_pS^2(a)$，平面$\operatorname{span}\{p,v\}$与球面相交为大圆，选适当方向和速度即有该初值；测地线初值唯一性说明所有非恒定球面测地线都在这种大圆上。常值曲线另计，不能称作一个大圆。

还可直接写出$\gamma(t)=\cos(|v|t/a)p+(a/|v|)\sin(|v|t/a)v$。它在球面上，初值为$p,v$，环境二阶导数为$-(|v|^2/a^2)\gamma$，与切空间正交；第11章切向投影联络等于Levi--Civita，所以其内蕴加速度为零。坐标极点处使用此全局公式，不把含$\cot\theta$的卡公式强行代入。
\end{proof}

\Needspace{12\baselineskip}
\subsection*{习题12.3：法坐标究竟拉直哪些线}
梁书定理3-3-8的文内主题改编。证明法坐标中从中心$p$出发的仿射测地线为$x(t)=tv$，并据此推出$\Gamma(p)=0$。这是否说明所有测地线都为坐标直线、整个邻域都平坦？
\begin{proof}[解答]
缩放初值的唯一性给$\gamma_v(t)=\exp_p(tv)$，因此用$\exp_p^{-1}$取坐标就得到$tv$。对任意$v$代入坐标测地线方程，在$t=0$有$\Gamma^k_{ij}(p)v^iv^j=0$；无挠使下槽对称，取$v=e_i$及$e_i+e_j$即可逐项得零。

只在中心得到系数为零。以球面法坐标为例，从中心出发的大圆成为径向直线，但第11章已独立算出$K=a^{-2}\ne0$，所以不可能在整个邻域$\Gamma=0$。不经过中心的测地线也没有在上述论证中被拉直。一般非坐标平行标架与法坐标的“中心零”亦不同，不能与第10章整个邻域的平坦性混淆。
\end{proof}

\Needspace{12\baselineskip}
\subsection*{习题12.4：下确界不达到的平坦例子}
原创。在$M=\R^2\setminus\{0\}$取欧氏度量，$p=(-1,0)$、$q=(1,0)$。求$d(p,q)$，证明没有最短曲线，并判断两种完备性。
\begin{proof}[解答]
任意路径的欧氏长度至少弦长2。对$0<\varepsilon<1$，沿$x$轴从$p$到$(-\varepsilon,0)$，绕上半圆半径$\varepsilon$，再沿轴到$q$，总长
\[
 L_\varepsilon=2-2\varepsilon+\pi\varepsilon
 =2+(\pi-2)\varepsilon\longrightarrow2.
\]
故$d(p,q)=2$。若达到等号，写速度为$(x',y')$，则$\int\sqrt{(x')^2+(y')^2}\geq\int x'=2$取等号，必有$y'=0$、$x'\geq0$几乎处处。连续性与分段光滑性使路径在$x$轴单调从$-1$到1，必经过删点，矛盾。

序列$(1/j,0)$沿正轴的距离为$|1/j-1/k|$，是Cauchy却在$M$无极限，所以距离不完备。测地线$\gamma(t)=(1-t,0)$在$t=1$撞删点，不能延拓，亦不测地完备。曲率恒零仍不排除这些现象；失败来自流形的缺口。
\end{proof}

\Needspace{12\baselineskip}
\subsection*{习题12.5：无共轭点也可能有多条最短线}
原创。对周长$2\pi a$的圆周，计算两点距离、指数映射、对跖点的最短线数及Jacobi场，解释它与指标引理不矛盾。
\begin{proof}[解答]
用局部弧长坐标$s$，度量为$\dd s^2$、联络系数为零。取初点弧长为0，切空间识别为$\R$，则$\exp_p(v)$为$v$模$2\pi a$的圆周点；若终点代表为$s$，
\[
 d(p,q)=\min_{k\in\mathbb Z}|s+2\pi ak|.
\]
为说明此式，对任意路径沿有限坐标弧逐段选连续角度，交叠只差$2\pi a$的整数倍；由起点值固定逐段拼成实线提升。提升终点为$s+2\pi ak$，长度等于提升在实线的长度，至少该绝对值；单调弧实现各候选，故取其最小值。

当$q$对跖于$p$，两个候选为$\pm\pi a$，有顺逆两条最短半圆；超过半圆的弧不再最短。可是维数1使$R=0$，平行单位场$T$下Jacobi场为$(A+Bt)T$，不能在两个不同时刻为零而非零。因此没有共轭点，指数微分处处非零。端点零场的$I=\int|V'|^2>0$与两条远处不同的最短线并不矛盾：正定指标型不提供全局路径唯一性。
\end{proof}

\Needspace{12\baselineskip}
\subsection*{习题12.6：三个常曲率Jacobi模型}
原创。在维数$n\geq2$的常截面曲率$K$模型中，沿单位速度测地线取平行单位法向场$E$。求初值
\[
 J(0)=0,\qquad D_tJ(0)=E(0)
\]
对应的Jacobi场。在球面$S^2(a)$据此求从一点出发的首次共轭距离和指数映射的角向微分。
\begin{proof}[解答]
第11章曲率模型给$R(E,T)T=KE$。写$J=s_K E$，得到$s_K''+Ks_K=0$、$s_K(0)=0,s_K'(0)=1$，故
\[
 s_K(t)=\begin{cases}
 \sin(\sqrt Kt)/\sqrt K,&K>0,\\
 t,&K=0,\\
 \sinh(\sqrt{-K}t)/\sqrt{-K},&K<0.
 \end{cases}
\]
正曲率时首次正零为$\pi/\sqrt K$；零与负曲率时无正零。球面$K=a^{-2}$给$s_K(t)=a\sin(t/a)$，所以首次共轭距离为$\pi a$，即对跖点，二维重数1。对$w\perp v$，$|v|=1$，命题\ref{prop:conjugate-exponential}给
\[
 (\dd\exp_p)_{tv}w=\frac{a\sin(t/a)}tP_{t\leftarrow0}w.
\]
径向单位输入仍被送为$T(t)$，所以只在角向退化；$t\to0$时系数趋1。该公式给局部微分的奇异性，不能仅凭它就推论所有全局最短性的分类。
\end{proof}

\Needspace{12\baselineskip}
\subsection*{习题12.7：球面第二变分的精确阈值}
原创。沿$S^2(a)$的一段单位速度大圆$\gamma:[0,b]\to S^2(a)$，取平行单位法向$E$与$V(t)=\sin(\pi t/b)E(t)$。计算$I(V,V)$，判断$b>\pi a$与$b=\pi a$的区别。
\begin{proof}[解答]
曲率项为$a^{-2}|V|^2$，所以
\[
 I(V,V)=\int_0^b\left(\frac{\pi^2}{b^2}\cos^2\frac{\pi t}b
 -\frac1{a^2}\sin^2\frac{\pi t}b\right)\dd t
 =\frac b2\left(\frac{\pi^2}{b^2}-\frac1{a^2}\right).
\]
若$b>\pi a$，指标负，第二变分给相同端点的更短曲线。若$b=\pi a$，指标为零，$V$与球面角向Jacobi场成比例；此时两端为对跖点，半大圆仍最短。为验证后一句而不预设球面距离公式，令$\theta$为点到起点径向的球心夹角。球面极坐标度量给任意路径$|\dot c|\geq a|\theta'|$；绕开坐标极点分段后取极限，起终夹角为0与$\pi$，所以长度至少$\pi a$，大圆半弧达到等号。

零二阶变化本身不能证明或否定最短；本例的最短性来自另一个长度下界。若取同一大圆更长一段，首次共轭点进入内部，便落入本章“不再最短”定理的范围。
\end{proof}

\Needspace{12\baselineskip}
\subsection*{习题12.8：开半平面为何可以完备}
原创。在$\mathbb H^2=\{(x,y):y>0\}$取$g=y^{-2}(\dd x^2+\dd y^2)$。证明其长度距离完备，推出测地完备，并验证$\gamma(t)=(0,e^t)$为单位速度测地线。
\begin{proof}[解答]
任意曲线满足$L\geq\int|y'|/y\geq|\log y_1-\log y_0|$。若$(x_j,y_j)$为$d$-Cauchy序列，则$\log y_j$为实数Cauchy序列，所以晚期$0<m\leq y_j\leq M<\infty$。

对晚期两点，取长度小于$\delta$的连接路径，其中$\delta$可随它们的距离趋零（用下确界取$L<d+\varepsilon$）。从起点到路径任一点仍有$|\log y-\log y_j|\leq\delta$，故沿途$y\leq Me^\delta$，从而
\[
 |x_j-x_k|\leq\int|x'|\leq Me^\delta\int\frac{|x'|}{y}
 \leq Me^\delta\delta.
\]
因此$x_j$亦为实数Cauchy序列，且极限点$(x,y)$满足$y>0$。在该点的紧坐标小球上度量与欧氏度量一致有界，命题\ref{prop:distance-topology}或直接直线长度上界说明$d((x_j,y_j),(x,y))\to0$。故完备，Hopf--Rinow给测地完备。

由第11章共形公式，非零系数为$\Gamma^x_{xy}=\Gamma^x_{yx}=-1/y$、$\Gamma^y_{xx}=1/y$、$\Gamma^y_{yy}=-1/y$。对$\gamma$，$x'=0,y'=y,y''=y$，故$y''-(y')^2/y=0$，另一方程也为零，且$|\dot\gamma|=y'/y=1$。它在$t\to-\infty$趋向边界，但需要无限弧长；“坐标域是欧氏开集”不能决定其黎曼完备性。
\end{proof}

